\def\topic{Category theory and the Yoneda lemma}
\ifdefined\MyHeaderIncluded
\else
\def\qwertyflag{true}

\ifdefined\MyHeaderIncluded

\else
  \newcommand{\foobar}{FooBar}
  \def\MyHeaderIncluded{true}
% if is till the end of the file

\documentclass[10pt,a4paper]{article}
\usepackage[utf8]{inputenc}
%\usepackage[german]{babel}
%\usepackage[cyrillic]{babel}
\usepackage[T1]{fontenc}
%\usepackage[T2A]{fontenc}
%\usepackage[english, russian]{babel}
\usepackage{amsmath}
\usepackage{amsfonts}
\usepackage{amssymb}
\usepackage{amsthm}
\usepackage{mathtools} %coloneqq

\usepackage{bussproofs} %

%\renewcommand{\thesubsection}{\arabic{subsection}}
%\newtheorem{thm}{Theorem}[section]
%\renewcommand{\thethm}{\arabic{subsection}.\arabic{thm}}

\usepackage{thmtools}
\usepackage{hyperref}
\hypersetup{
    colorlinks=true,
    linkcolor=blue,
    filecolor=magenta,      
    urlcolor=cyan,
    pdftitle={Overleaf Example},
    %pdfpagemode=FullScreen,
    }
%\usepackage{stix-type1}%\turnediota
\theoremstyle{definition}

\newtheorem{constructor}{Constructor}[section]
\newtheorem{iaxiom}{Axiom}
\newtheorem{axiom}{Axiom}
\newtheorem*{axiom*}{Axiom}
\newtheorem{definition}{Definition}[section]
\newtheorem*{definition*}{Definition}
\newtheorem{remark}{Remark}[section]
\newtheorem*{remark*}{Remark}
\newtheorem{recall}{Recall}[section]
\newtheorem*{recall*}{Recall}

\newtheorem{lemma}{Lemma}[section]
\newtheorem{theorem}{Theorem}[section]
\newtheorem*{theorem*}{Theorem}
\newtheorem{note}{Note}[section]
\usepackage{etoolbox}

%\newcommand*{\@Dashv}[2]{%
%  \reflectbox{$\m@th#1#2$}%
%}
\newcommand*{\separation}[3]{% 1)binder, 2)set, 3)condition
  \{#1\in #2: #3\}
}

% https://www.freecodecamp.org/news/simple-regex-tricks-for-beginners-3acb3fa257cb/
%  \\mathcal\{T\} |---> \topology
%  \\mathcal\{([^\}]*)\}   |--->  $1
\newcommand{\AND}{\ \ \&\ \ }
\newcommand{\any}{\mathrm{\ any}}
\newcommand{\borel}{\mathcal{B}}
\newcommand{\can}{\mathrm{can}}
\newcommand{\canonize}{\mathrm{canonize}}
\newcommand{\closure}{\mbox{cl}}
\newcommand{\complex}{\mathbb{C}}
\newcommand{\converges}{\longrightarrow}
\newcommand{\countable}{\mbox{\ countable}}
\newcommand{\ldefined}{\mbox{\ defined}}
\newcommand{\dom}{\mbox{dom}}
\newcommand{\Dom}{\mbox{Dom}}
\newcommand{\Done}{\noindent \textbf{Done.}\\}
\newcommand{\lexists}{\mbox{\ exists}}
\newcommand{\lforall}{\mbox{\ forall}}
\newcommand{\lunique}{\mbox{\ unique}}
\newcommand{\family}{\mathrm{\ family}}
\newcommand{\Fm}{\mathrm{Fm}}
\newcommand{\FuncSymb}{\mathrm{FuncSymb}}
\newcommand{\function}{\mathrm{\ function}}
\newcommand{\functional}{\mbox{\ functional}}
\newcommand{\IFF}{\ \Longleftrightarrow\ }
\renewcommand{\iff}{\ \longleftrightarrow\ }
\newcommand{\IMPLIES}{\ \Longrightarrow\ }
\renewcommand{\implies}{\ \longrightarrow\ }
\newcommand{\Inf}{\mathrm{Inf}}
\newcommand{\injective}{\mbox{\ injective}}
\newcommand{\injection}{\mbox{\ injection}}
\newcommand{\interior}{\mbox{int}}
\newcommand{\intersection}{\bigcap}
\newcommand{\isempty}{\mathrm{\ empty}}
\newcommand{\LB}{\mathrm{LB}}
\newcommand{\nghbset}{\mathcal{N}}
\newcommand{\nonempty}{\mathrm{\ nonempty}}
\newcommand{\Ord}{\mathrm{Ord}}
\newcommand{\power}{\mathcal{P}}
\newcommand{\pr}{\mathrm{pr}}
\newcommand{\fst}{\mathrm{fst}}
\newcommand{\snd}{\mathrm{snd}}
\newcommand{\PredSymb}{\mathrm{PredSymb}}
\newcommand{\Proof}{\noindent \textbf{Proof.}\\}
\newcommand{\ran}{\mathrm{ran}}
\newcommand{\Ran}{\mathrm{Ran}}
\newcommand{\real}{\mathbb{R}}
\newcommand{\relation}{\mathrm{\ relation}}
\newcommand{\set}{\mathrm{\ set}}
\newcommand{\singleton}{\mathrm{\ singleton}}
\newcommand{\semantics}{\mathrm{semantics}}
\newcommand{\semialgebra}{\mathcal{S}}
\newcommand{\Sup}{\mathrm{Sup}}
\newcommand{\surjective}{\mbox{\ surjective\ on\ }}
\newcommand{\Term}{\mathrm{Term}}
\newcommand{\Type}{\mathrm{Type}}
\newcommand{\topology}{\mathcal{T}}
\newcommand{\total}{\mbox{\ total\ on\ }}
\newcommand{\transitive}{\mathrm{\ transitive}}
\newcommand{\union}{\bigcup}
\newcommand{\lif}{\mathrm{if}}
\newcommand{\lthen}{\mathrm{then}}
\newcommand{\lelse}{\mathrm{else}}
\newcommand{\lexception}{\mathrm{exception}}
\newcommand{\lundef}{\mathrm{undef}}
\newcommand{\univocal}{\mathrm{\ univocal}}
\newcommand{\univocalat}{\mathrm{\ univocal\ at\ }}
\newcommand{\UB}{\mathrm{UB}}
\newcommand{\Var}{\mathrm{Var}}

\newcommand{\notvarphi}{\phi}
\renewcommand{\phi}{\varphi}

\newcommand{\comment}[1]{\quad\mbox{(#1)}}
% local commands
%\preto\section{\undef\blipo} %etoolbox
\preto\subsection{\undef\blipo} %etoolbox
% https://tex.stackexchange.com/questions/105684/limit-macro-newcommand-scope-to-section
\usepackage{graphicx}
\usepackage{authblk}

\setlength\parindent{0pt} % whole doc of \noindent

%\usepackage{mathgl}

\fi % guard
\ifdefined\topic
\else
\def\topic{No Titel}
\fi
\title{\topic}
\author[1]{Georgy Dunaev}
\begin{document}
\maketitle
\fi
\section{Category theory}

This section develops the author's Hom-functor construction and its
Yoneda correspondence. The essential distinction is between
\emph{precomposition}, which changes the source of a morphism, and
\emph{postcomposition}, which changes its target.

We write $C$ for a category and, when the context is clear, also for its
collection of objects $\Ob_C$. Composition is read from right to left:
if $g:X\to Y$ and $f:Y\to Z$, then $f\circ g:X\to Z$. Thus its typed form is
\[
\circ:\Hom_C(Y,Z)\times\Hom_C(X,Y)\longrightarrow\Hom_C(X,Z).
\]
The opposite category has the same objects and
$\Hom_{C^\op}(X,Y)=\Hom_C(Y,X)$, with composition reversed.

We assume that $C$ is locally small: each $\Hom_C(X,Y)$ is a set.
The notation $\Set$ denotes the category of sets in the chosen size
convention. Functors and natural transformations are considered in an
ambient framework large enough to contain them; local smallness alone
does not make the collection of all functors a set. One may, for example,
work with a small $C$ relative to a chosen universe and use a larger
ambient universe for the functor category. The calculations below are
componentwise and do not depend on this particular choice of framework.

\begin{note}[Comparison with the Obelisk Yoneda summary]
The newer portal's \texttt{yoneda\_ai.html} gives the same correspondence
$\eta\mapsto\eta_A(1_A)$ and inverse $u\mapsto(f\mapsto G(f)(u))$
proved below. Its discussion of ``functor classes'' needs a size
interpretation: an individual functor may have a class-sized graph,
while a collection of all such graphs can lie at a higher level.
Impredicative class comprehension still has sets as members, so it does
not by itself turn proper-class graphs into members of a class.
The ambient convention stated above makes this distinction explicit
without changing the componentwise Yoneda argument.
\end{note}

\subsection{Functors and the contravariant Hom-functor}
\begin{definition}[Functor]
A covariant functor from $C$ to $D$ consists of an object map $F$ and
maps on morphisms
\[
\FMor^{A,B}:\Hom_C(A,B)\longrightarrow\Hom_D(F(A),F(B))
\]
that preserve identities and composition:
\[
\FMor^{A,A}(1_A)=1_{F(A)},\qquad
\FMor(f\circ g)=\FMor(f)\circ\FMor(g).
\]
When no ambiguity arises, we write $F(f)$ for $\FMor(f)$.
\end{definition}

\begin{definition}[Functor category]
The notation $[C,D]$ denotes the category whose objects are functors
$C\to D$ and whose morphisms are natural transformations. A natural
transformation $\eta:F\to G$ consists of components
$\eta_A:F(A)\to G(A)$ satisfying
\[
G(f)\circ\eta_A=\eta_B\circ F(f)
\quad\text{for every }f:A\to B.
\]
\end{definition}

\begin{definition}[Morphism action of the Hom-functor]
Fix $x\in C$. On objects, put $h_x(A)=\Hom_C(A,x)$. For a morphism
$f:A\to B$ of $C$, define precomposition by
\[
\HomMor_C^{A,B}(f,x):\Hom_C(B,x)\longrightarrow\Hom_C(A,x),
\qquad g\longmapsto g\circ f.
\]
This is the morphism action of $h_x$ on the corresponding arrow
$B\to A$ of $C^\op$. The notations $\Hom_C(-,x)$ and $h_x$ denote the
same functor; $\HomMor_C(f,x)$ makes its morphism action explicit.
\end{definition}

\begin{lemma}[Contravariant Hom-functor]
For each $x\in C$, the construction above gives
$h_x=\Hom_C(-,x)\in[C^\op,\Set]$.
\end{lemma}
\Proof
The object values are sets by local smallness. Precomposition has the
displayed source and target, because composition in $C$ is defined for
these morphisms. It preserves identities since $g\circ1_A=g$.
For $f:A\to B$, $k:B\to D$, and $g:D\to x$, associativity gives
\[
\HomMor_C(k\circ f,x)(g)
 =g\circ(k\circ f)
 =(g\circ k)\circ f
 =\bigl(\HomMor_C(f,x)\circ\HomMor_C(k,x)\bigr)(g).
\]
This reversed order is exactly the composition law for a functor on
$C^\op$. Equality on each argument gives equality of the corresponding
maps.
\Done

\subsection{The Yoneda functor}
\begin{definition}[Yoneda functor]
On objects, the Yoneda functor sends $x$ to $h_x=\Hom_C(-,x)$.
For $f:x\to y$, define $h_f:h_x\to h_y$ by postcomposition:
\[
(h_f)_w:\Hom_C(w,x)\longrightarrow\Hom_C(w,y),
\qquad g\longmapsto f\circ g.
\]
This component family is the transformation denoted by $\eta$ in the
construction below.
\end{definition}

\begin{lemma}[Naturality of postcomposition]
For every $f:x\to y$, the family $\eta=h_f$ is a natural transformation.
\end{lemma}
\Proof
An arrow $m:a\to b$ in $C^\op$ is an arrow $m:b\to a$ in $C$.
For $g:a\to x$, the two sides of the naturality square give
\[
\bigl(\eta_b\circ\HomMor_C(m,x)\bigr)(g)=f\circ(g\circ m),
\]
\[
\bigl(\HomMor_C(m,y)\circ\eta_a\bigr)(g)=(f\circ g)\circ m.
\]
Associativity makes them equal. This also identifies the source and
target of both maps as $\Hom_C(a,x)\to\Hom_C(b,y)$.
\Done

\begin{lemma}[Functoriality of the Yoneda construction]
The assignments $x\mapsto h_x$ and $f\mapsto h_f$ define a functor
$h_-:C\to[C^\op,\Set]$.
\end{lemma}
\Proof
For every $g:w\to x$, $(h_{1_x})_w(g)=1_x\circ g=g$, so identities
are preserved. If $f:x\to y$ and $k:y\to z$, then
\[
(h_{k\circ f})_w(g)=(k\circ f)\circ g
 =k\circ(f\circ g)
 =\bigl((h_k)_w\circ(h_f)_w\bigr)(g).
\]
Natural transformations compose componentwise, so composition is
preserved as well.
\Done

\subsection{The Yoneda correspondence}
\begin{lemma}[Yoneda lemma, contravariant version]
Let $G:C^\op\to\Set$ and $A\in C$. There is a bijection
\[
\Nat(h_A,G)\longrightarrow G(A),\qquad
\Psi\longmapsto\Psi_A(1_A).
\]
Its inverse sends $u\in G(A)$ to the transformation $\Phi^u$ with
\[
\Phi^u_x(f)=G(f)(u)\qquad(f:x\to A\text{ in }C).
\]
Here $G(f):G(A)\to G(x)$ denotes the action on the corresponding
arrow of $C^\op$.
\end{lemma}
\Proof
First verify that the inverse formula defines a natural transformation.
For $m:y\to x$ and $f:x\to A$ in $C$, functoriality gives
\[
G(m)\bigl(\Phi^u_x(f)\bigr)
 =G(m)\bigl(G(f)(u)\bigr)
 =G(f\circ m)(u)
 =\Phi^u_y(f\circ m).
\]
This is the naturality condition for $\Phi^u:h_A\to G$.
Its component at the identity recovers the chosen element:
\[
\Phi^u_A(1_A)=G(1_A)(u)=u.
\]
Thus different choices of $u$ give different transformations.

Conversely, let $\Psi:h_A\to G$ be natural and put
$u=\Psi_A(1_A)$. For any $f:x\to A$, naturality on the corresponding
arrow $A\to x$ in $C^\op$ gives
\[
\Psi_x(f)
 =\Psi_x(1_A\circ f)
 =G(f)\bigl(\Psi_A(1_A)\bigr)
 =G(f)(u)
 =\Phi^u_x(f).
\]
The equality holds for every $x$ and every $f$, so $\Psi=\Phi^u$.
The two displayed constructions are therefore inverse bijections.
\Done

The correspondence is natural in both parameters. A transformation
$\alpha:G\to H$ sends $\Psi_A(1_A)$ to
$\alpha_A(\Psi_A(1_A))$, which is the element corresponding to
$\alpha\circ\Psi$. For an arrow $k:A\to B$, precomposing
$\Psi:h_B\to G$ with $h_k$ gives the element
\[
(\Psi\circ h_k)_A(1_A)=\Psi_A(k)
 =G(k)\bigl(\Psi_B(1_B)\bigr).
\]
These identities state precisely how the bijection respects a change of
functor or representing object.

\begin{lemma}[Yoneda embedding]
The functor $h_-:C\to[C^\op,\Set]$ is fully faithful.
\end{lemma}
\Proof
Apply the Yoneda correspondence with $A=x$ and $G=h_y$. It gives
\[
\Nat(h_x,h_y)\longrightarrow\Hom_C(x,y),\qquad
\eta\longmapsto\eta_x(1_x).
\]
The inverse sends $f:x\to y$ to the transformation with components
$g\mapsto f\circ g$, exactly $h_f$. Hence the map induced by $h_-$
on each Hom-set is a bijection: it is surjective (full) and injective
(faithful).
\Done

\subsection{Universal elements and uniqueness of representations}
A representation is more than an object whose elements have the desired
shape: it includes a natural identification. The following formulation
makes that extra structure explicit.

\begin{definition}[Universal element]
For $G:C^\op\to\Set$, an element $u\in G(A)$ is universal if, for every
object $X$ and $v\in G(X)$, there is exactly one morphism $f:X\to A$
with $G(f)(u)=v$.
\end{definition}

\begin{proposition}[Universal-element criterion]
The functor $G$ is naturally isomorphic to $h_A$ if and only if it has a
universal element in $G(A)$. If $(A,u)$ and $(B,v)$ are two such choices,
there is a unique isomorphism $k:A\to B$ satisfying $G(k)(v)=u$.
\end{proposition}
\Proof
The Yoneda transformation $\Phi^u_X:f\mapsto G(f)(u)$ is natural.
Universality says exactly that every component is bijective. The inverses
are natural: from
$G(m)\Phi^u_X=\Phi^u_Yh_A(m)$ for $m:Y\to X$, compose with the inverse
components to obtain
$(\Phi^u_Y)^{-1}G(m)=h_A(m)(\Phi^u_X)^{-1}$.
Conversely, for a natural isomorphism $\Phi:h_A\to G$, put
$u=\Phi_A(1_A)$. Yoneda gives $\Phi=\Phi^u$, and componentwise bijectivity
gives universality.

For two universal elements there are unique arrows $k:A\to B$ and
$l:B\to A$ with $G(k)(v)=u$ and $G(l)(u)=v$. Contravariance gives
\[
G(l\circ k)(u)=G(k)G(l)(u)=u,
\qquad
G(k\circ l)(v)=G(l)G(k)(v)=v.
\]
Uniqueness for $(A,u)$ and $(B,v)$ respectively implies
$l\circ k=1_A$ and $k\circ l=1_B$. The condition on $k$ already determines
it uniquely.
\Done

\begin{example}[Products as representing objects]
Fix objects $B,D$ in $C$ and define
$G(X)=\Hom_C(X,B)\times\Hom_C(X,D)$, with
$G(m)(b,d)=(b\circ m,d\circ m)$. Identity and composition laws follow
componentwise from those in $C$. An element of $G(P)$ is a pair of arrows
$p:P\to B$, $q:P\to D$. Its universality says exactly that every pair
$b:X\to B$, $d:X\to D$ factors uniquely as
\[
p\circ f=b,\qquad q\circ f=d\qquad(f:X\to P).
\]
Thus a categorical product represents this functor. The preceding
proposition proves uniqueness up to the unique isomorphism preserving
both projections. It does not assert that products exist in every category.
\end{example}

\begin{proposition}[Yoneda detects isomorphisms]
An arrow $f:A\to B$ is an isomorphism if and only if every map
\[
\Hom_C(X,A)\longrightarrow\Hom_C(X,B),\qquad g\longmapsto f\circ g
\]
is bijective.
\end{proposition}
\Proof
An inverse to $f$ supplies the inverse maps by postcomposition.
Conversely, surjectivity at $X=B$ gives $g:B\to A$ with $f\circ g=1_B$.
Then $f\circ(g\circ f)=f=f\circ1_A$. Injectivity at $X=A$ gives
$g\circ f=1_A$, so $g$ is an inverse.
\Done

The quantifier over all test objects is essential. In the category of
abelian groups, the zero group is terminal, and there is just one arrow
from it to any group. Testing only that object therefore makes every
arrow appear bijective on these Hom-sets, including the noninvertible
arrow from the zero group to $\mathbb Z$. In $\Set$, by contrast,
maps from a singleton identify with elements, and this particular test
object does detect bijections.


\subsection{Cancellation and pointwise tests}
\begin{proposition}[What Hom-functors detect about cancellation]
An arrow $f:A\to B$ is a monomorphism if and only if postcomposition
$\Hom_C(X,A)\to\Hom_C(X,B)$ is injective for every $X$.
It is an epimorphism if and only if precomposition
$\Hom_C(B,X)\to\Hom_C(A,X)$ is injective for every $X$.
By contrast, postcomposition is surjective for every $X$ if and only
if $f$ has a right inverse.
\end{proposition}
\Proof
The first two equivalences are precisely left and right cancellation,
respectively: $f\circ g=f\circ h$ implies $g=h$, or
$g\circ f=h\circ f$ implies $g=h$, for each permissible source or
target. If postcomposition is surjective at $X=B$, the identity of $B$
has a preimage $s:B\to A$, so $f\circ s=1_B$. Conversely, such an
$s$ lifts every $g:X\to B$ to $s\circ g:X\to A$, proving surjectivity
at all $X$.
\Done

\begin{example}[An epimorphism need not give pointwise surjectivity]
In the category of abelian groups, the quotient
$q:\mathbb Z\to\mathbb Z/2\mathbb Z$ is an epimorphism: two
homomorphisms agreeing after $q$ agree on every element, because $q$
is surjective as a set map. But it has no right inverse. Any
homomorphism $s:\mathbb Z/2\mathbb Z\to\mathbb Z$ sends the generator
to an integer $n$ satisfying $2n=0$, so $n=0$ and $s=0$.
Thus postcomposition with $q$ is not surjective on the Hom-set with
source $\mathbb Z/2\mathbb Z$. This distinguishes cancellation in the
category from pointwise surjectivity of the corresponding natural
transformation $h_q$; the latter is the stronger split-epimorphism
condition proved above.
\end{example}


\subsection{Hom-functors and universal constructions}
There are two variances to keep apart. For a fixed source $X$, the
\emph{covariant} functor $\Hom_C(X,-)$ sends a target arrow to
postcomposition. For a fixed target $A$, the earlier $h_A=\Hom_C(-,A)$
is contravariant. The next statement concerns the former functor.

\begin{proposition}[Products and equalizers tested by Hom]
If $P$ with projections $p:P\to B$, $q:P\to D$ is a product, then
for every $X$ the map
\[
\Hom_C(X,P)\longrightarrow\Hom_C(X,B)\times\Hom_C(X,D),
\qquad u\longmapsto(p\circ u,q\circ u)
\]
is bijective. If $e:E\to B$ is an equalizer of $f,g:B\to D$, then
postcomposition with $e$ gives a bijection
\[
\Hom_C(X,E)\longrightarrow
\{u\in\Hom_C(X,B):f\circ u=g\circ u\}.
\]
Both bijections are natural in the test object $X$.
\end{proposition}
\Proof
The product property gives existence and uniqueness of $u$ with any
prescribed two components, which proves the first bijection. The
equalizer property says exactly that every arrow $u$ satisfying the
displayed equality has a unique factorization $u=e\circ v$, proving
the second. For $k:Y\to X$, precomposing $u$ or $v$ by $k$ precomposes
all displayed components by $k$, by associativity. This proves naturality.
\Done

These statements prove preservation of the indicated limits by
$\Hom_C(X,-)$; they do not assume those limits exist in every category.
Dually, the contravariant $h_A$ takes an existing coproduct to a product
of Hom-sets and an existing coequalizer to an equalizer of Hom-sets.
The proof is the same factorization argument with arrows reversed.
All Hom-sets here are sets by local smallness; no collection of all
proper-class-valued diagrams is being formed.

\subsection{Small limits from products and equalizers}
A \emph{small diagram} here is a functor $D:J\to C$, where the
objects and arrows of the indexing category $J$ form sets in the chosen
size convention. No smallness of all of $C$ is assumed. A cone from
$X$ to $D$ is a family $u_j:X\to D(j)$ satisfying
\[
D(\alpha)\circ u_i=u_j\qquad(\alpha:i\to j\text{ in }J).
\]
A limit is a cone $(L,(\ell_j))$ such that each cone $(X,(u_j))$
has a unique factorization $u_j=\ell_j\circ h$ for all $j$, with
$h:X\to L$. For a set-indexed family, a product has the same
factorization property without compatibility equations between its
components; an empty product is therefore a terminal object.

\begin{theorem}[Construction of a small limit]
\label{thm:small_limits_products_equalizer}
For the specified small diagram $D$, suppose the two products
\[
A=\prod_{j\in\Ob_J}D(j),\qquad
B=\prod_{\alpha:i\to j\text{ in }J}D(j)
\]
exist in $C$. Let $p_j:A\to D(j)$ be the first product's projections,
and let $q_\alpha:B\to D(j)$ be the second product's projections.
There are unique arrows $s,t:A\to B$ determined by
\[
q_\alpha\circ s=D(\alpha)\circ p_i,\qquad
q_\alpha\circ t=p_j\qquad(\alpha:i\to j).
\]
If an equalizer $e:L\to A$ of $s,t$ exists, then
$\ell_j=p_j\circ e$ makes $L$ a limit of $D$.
Consequently a category with all small products and equalizers has
limits of every small diagram.
\end{theorem}
\Proof
The product property of $B$ gives $s,t$ from the displayed families
of arrows with common source $A$. Since $s\circ e=t\circ e$,
for every $\alpha:i\to j$ we have
\[
D(\alpha)\circ\ell_i
=q_\alpha\circ s\circ e
=q_\alpha\circ t\circ e=\ell_j.
\]
Thus the proposed family is a cone.

Given any cone $(X,(u_j))$, the product property of $A$ gives a
unique $u:X\to A$ with $p_j\circ u=u_j$ for all $j$.
For each $\alpha:i\to j$,
\[
q_\alpha\circ s\circ u=D(\alpha)\circ u_i
=u_j=q_\alpha\circ t\circ u.
\]
Uniqueness in the product property of $B$ gives $s\circ u=t\circ u$.
The equalizer property therefore gives a unique $h:X\to L$ with
$e\circ h=u$, and hence $\ell_j\circ h=u_j$.
If $h':X\to L$ has the same cone components, then
$p_j\circ e\circ h'=u_j$ for all $j$. Product uniqueness gives
$e\circ h'=u$, and equalizer uniqueness gives $h'=h$.
This proves the limit property.

If $J$ has no objects, it has no arrows. Both products are terminal
objects, and $s=t$ is the unique arrow between them. One may take
$e=1_A$ as their equalizer, giving the terminal object $A$ as the
limit of the empty diagram. More generally the preceding argument
shows that any chosen equalizer in this case is terminal.
The proof constructs the limit for one specified diagram from specified
products and an equalizer; it requires no simultaneous choice of limits
for a class of diagrams.
\Done

\begin{proposition}[Hom preserves and detects the limit property]
\label{prop:hom_small_limits}
A cone $(L,(\ell_j))$ to a small diagram $D:J\to C$ is a limit
if and only if, for every $X\in C$, the map
\[
\begin{split}
\Hom_C(X,L)&\longrightarrow\operatorname{Cone}(X,D),\\
h&\longmapsto(\ell_j\circ h)_{j\in\Ob_J}
\end{split}
\]
is bijective. The set on the right consists of the compatible families
in the product
\[
\prod_{j\in\Ob_J}\Hom_C(X,D(j)).
\]
These bijections are natural in $X$. In particular, the covariant
functor $\Hom_C(X,-)$ preserves every existing small limit.
\end{proposition}
\Proof
Local smallness and the set-sized index make the displayed product
of Hom-sets, and its subset of compatible families, sets in the stated
ambient convention. A cone determines the displayed map because
precomposing its compatibility equations by $h$ gives those for the
resulting family. Surjectivity is existence of a factorization for
each cone, and injectivity is its uniqueness. This proves the first
assertion in both directions.

The compatible-family set is the limit in $\Set$ of the diagram
$j\mapsto\Hom_C(X,D(j))$, whose arrow action is postcomposition by
$D(\alpha)$. Indeed, the coordinate projections form a cone; any
compatible family of functions from a set $S$ determines a unique
function from $S$ into the compatible-family set by evaluation at each
$s\in S$. Thus the preceding bijection says precisely that
$\Hom_C(X,-)$ preserves the given limit.
For an arrow $k:Y\to X$, both routes send $h:X\to L$ to the
family $(\ell_j\circ h\circ k)_j$, by associativity. This proves
naturality in the test object. In the empty-diagram case the
compatible-family set is the singleton containing the empty family,
so the criterion reduces to having exactly one arrow $X\to L$
from each $X$, as required for a terminal object.
\Done

The indexing restriction is essential to the size claim: local
smallness controls individual Hom-sets, not arbitrary proper-class
products of them. These theorems assert the stated universal properties
inside the chosen framework and do not construct limits of every
class-indexed diagram.

\subsection{Initial algebras, final coalgebras, and fixed points}
Let $P:C\to C$ be an endofunctor. A $P$-algebra is an arrow
$a:P(A)\to A$. An algebra morphism from $(A,a)$ to $(B,b)$ is an
arrow $u:A\to B$ with
\[
u\circ a=b\circ P(u).
\]
Identities satisfy the equation because $P(1_A)=1_{P(A)}$.
For algebra morphisms $u:(A,a)\to(B,b)$ and $v:(B,b)\to(D,c)$,
\[
(v\circ u)\circ a
=v\circ b\circ P(u)
=c\circ P(v)\circ P(u)
=c\circ P(v\circ u).
\]
Thus composition is inherited from $C$, and its associativity and identity
laws are inherited as well. These objects and arrows therefore form an
algebra category in the chosen ambient size convention.
An algebra is \emph{initial} when it has exactly one
algebra morphism to every algebra. Dually, a coalgebra is an arrow
$d:M\to P(M)$, a coalgebra morphism $u:(M,d)\to(N,e)$ satisfies
$e\circ u=P(u)\circ d$, and a coalgebra is \emph{final} when every
coalgebra has exactly one morphism to it. The analogous composition
calculation, with the structure arrows reversed, verifies the coalgebra
category laws.

\begin{lemma}[Lambek's lemma]
\label{lem:lambek}
The structure map of an initial algebra is an isomorphism. The
structure map of a final coalgebra is also an isomorphism.
\end{lemma}
\Proof
Let $(A,a)$ be initial. The arrow $P(a):P(P(A))\to P(A)$ is another
algebra, so there is a unique algebra morphism
$h:(A,a)\to(P(A),P(a))$. The arrow $a$ itself is an algebra morphism
in the reverse direction: its required equation is the identity
$a\circ P(a)=a\circ P(a)$. Thus $a\circ h$ is an endomorphism of
the initial algebra, and uniqueness gives $a\circ h=1_A$.
The equation for $h$ now gives
\[
h\circ a=P(a)\circ P(h)=P(a\circ h)=1_{P(A)}.
\]
Hence $h$ is an inverse to $a$.

For a final coalgebra $(M,d)$, let
$k:(P(M),P(d))\to(M,d)$ be its unique coalgebra morphism.
The arrow $d$ is a coalgebra morphism in the other direction, so finality
gives $k\circ d=1_M$. Its morphism equation yields
\[
d\circ k=P(k)\circ P(d)=P(k\circ d)=1_{P(M)}.
\]
Thus $k$ is an inverse to $d$.
\Done

\begin{proposition}[Uniqueness with structure]
Two initial $P$-algebras are uniquely isomorphic as algebras, and two
final $P$-coalgebras are uniquely isomorphic as coalgebras.
\end{proposition}
\Proof
For initial algebras, take their unique morphisms in the two directions.
Each composite is an endomorphism of an initial object, hence equals
its identity. The two morphisms are therefore inverse, and uniqueness
is part of initiality. Reverse the arrows for final coalgebras.
\Done

For the polynomial functor
$P_B(X)=\sum_{a\in A}X^{B(a)}$ on sets, the W-type construction in
Theorem~\ref{thm:general_W_recursion} supplies an initial algebra, and
Theorem~\ref{thm:M_final_coalgebra} supplies the finite-address M-type
as a final coalgebra.
Lambek's lemma explains their fixed-point equations without treating a
bare equation $P_B(X)\cong X$ as a substitute for the universal property.
The recursion and corecursion theorems supply the unique structure-preserving
maps; the preceding proposition then makes the constructions independent
of their particular encodings, up to those unique isomorphisms.

\subsection{Fusion, induction, and coinduction}
Keep the same endofunctor $P:C\to C$. For an initial algebra $(A,a)$,
write $\operatorname{fold}_b:A\to B$ for the unique morphism to an
algebra $(B,b)$. For a final coalgebra $(M,d)$, write
$\operatorname{unfold}_c:X\to M$ for the unique morphism from a
coalgebra $(X,c)$. These names refer to the universal properties,
not to an additional algorithm or a choice of encodings.

\begin{proposition}[Recursion and corecursion fusion]
If $h:(B,b)\to(D,e)$ is an algebra morphism, then
\[
h\circ\operatorname{fold}_b=\operatorname{fold}_e.
\]
If $h:(X,c)\to(Y,e)$ is a coalgebra morphism, then
\[
\operatorname{unfold}_e\circ h=\operatorname{unfold}_c.
\]
\end{proposition}
\Proof
For the algebra assertion, the defining equations and functoriality give
\[
\begin{aligned}
(h\circ\operatorname{fold}_b)\circ a
 &=h\circ b\circ P(\operatorname{fold}_b)\\
 &=e\circ P(h)\circ P(\operatorname{fold}_b)\\
 &=e\circ P(h\circ\operatorname{fold}_b).
\end{aligned}
\]
Thus the composite is an algebra morphism from the initial algebra to
$(D,e)$, so uniqueness identifies it with $\operatorname{fold}_e$.
For the coalgebra assertion, similarly,
\[
\begin{aligned}
d\circ(\operatorname{unfold}_e\circ h)
 &=P(\operatorname{unfold}_e)\circ e\circ h\\
 &=P(\operatorname{unfold}_e)\circ P(h)\circ c\\
 &=P(\operatorname{unfold}_e\circ h)\circ c.
\end{aligned}
\]
The composite is therefore a coalgebra morphism to the final coalgebra,
and uniqueness identifies it with $\operatorname{unfold}_c$.
\Done

For the polynomial W-types, the first identity says that transforming
a recursively computed result by a map preserving the constructor
operations gives the same result as the corresponding recursion into
the new target. For the M-types, the second says that a map of state
systems preserving their labelled transitions does not change their
unfolded trees. The preservation equations for $h$ are essential;
arbitrary maps need not satisfy these identities.

\begin{proposition}[Subalgebra and quotient rigidity]
Suppose $(A,a)$ is initial and an algebra morphism
$m:(S,s)\to(A,a)$ has an underlying arrow $m:S\to A$ that is a
monomorphism in $C$. Then $m$ is an isomorphism of algebras.
Dually, suppose $(M,d)$ is final and a coalgebra morphism
$q:(M,d)\to(Q,e)$ has an underlying arrow $q:M\to Q$ that is an
epimorphism in $C$. Then $q$ is an isomorphism of coalgebras.
\end{proposition}
\Proof
Initiality supplies an algebra morphism $h:(A,a)\to(S,s)$.
The composite $m\circ h$ is an endomorphism of the initial algebra,
so it equals $1_A$. Consequently
\[
m\circ(h\circ m)=(m\circ h)\circ m=m=m\circ1_S.
\]
Left cancellation by $m$ gives $h\circ m=1_S$. Thus $m$ and $h$
are inverse algebra morphisms.

For the dual assertion, finality supplies a coalgebra morphism
$k:(Q,e)\to(M,d)$, and $k\circ q=1_M$ by uniqueness of the
endomorphism of the final coalgebra. Hence
\[
(q\circ k)\circ q=q\circ(k\circ q)=q=1_Q\circ q.
\]
Right cancellation by $q$ gives $q\circ k=1_Q$, so $q$ and $k$
are inverse coalgebra morphisms.
\Done

Both assertions assume that the displayed structure on the subobject
or quotient already exists and makes the indicated arrow a morphism.
They do not claim that every subobject is closed under the operations,
or that every quotient admits a compatible coalgebra structure. No
preservation of monomorphisms or epimorphisms by $P$ is assumed.

In sets, this gives a direct form of structural induction for the
W-type in Theorem~\ref{thm:general_W_recursion}. Let $S\subseteq W$
be a subset such that grafting any labelled family of children in $S$
again gives an element of $S$. Restricting grafting defines
$s:P_B(S)\to S$, and the inclusion $m:S\hookrightarrow W$ satisfies
$m\circ s=a\circ P_B(m)$. It is a monomorphism, so the proposition
makes it an isomorphism. Since it is the inclusion, this says $S=W$.
This derives induction from initiality for predicates whose extension
is a specified subset, independently of a particular tree encoding.

There is also a useful direct coinduction consequence of finality.
If $(T,c)$ is a coalgebra and $u,v:(T,c)\to(M,d)$ are coalgebra
morphisms, then $u=v$, since the target is final. In particular, let
$R\subseteq M_B\times M_B$ satisfy the polynomial bisimulation
condition: for each $(t,u)\in R$, the root labels agree, say at $a$,
and $(t_b,u_b)\in R$ for every $b\in B(a)$. Then
\[
(t,u)\longmapsto\bigl(a,\ b\longmapsto(t_b,u_b)\bigr)
\]
defines a coalgebra structure on the set $R$. Both projections from
$R$ to $M_B$ are coalgebra morphisms, by the displayed matching of
labels and children. Finality in Theorem~\ref{thm:M_final_coalgebra}
makes the projections equal, so every related pair has $t=u$.
The structure on $R$ is explicitly constructed; no arbitrary lifting
of relations through an endofunctor or choice of matching successors
is being assumed. This recovers the tree coinduction principle from
its universal property.

\begin{example}[A fixed-point equation alone does not give initiality]
Take $P$ to be the identity functor on sets. The map $1_X:X\to X$ is
an isomorphism for every set $X$, so every such $X$ satisfies the bare
fixed-point equation. But the singleton algebra $(\{0\},1_{\{0\}})$
is not initial: there is no map from it to the empty algebra.
Dually, the empty coalgebra is not final because the singleton coalgebra
has no map into it. For this functor the empty identity algebra is
initial, and the singleton identity coalgebra is final. Thus even an
invertible structure map is strictly less information than the relevant
universal property.
\end{example}

\subsection{Locally Small Categories, the Yoneda Lemma, and Pointwise Kan Extensions}
\label{subsec:yoneda_presheaves_and_kan_extensions}

In \texttt{MKTD1-theory} and \texttt{MKTD2-theory}, a locally small category $\mathcal{C}$ with a proper class of objects $\operatorname{Ob}(\mathcal{C}) \subseteq \V$ represents each hom-set $\operatorname{Hom}_{\mathcal{C}}(X, Y) \in \V$ via Dunaev Tagged Ordered Pairs $\pair{\pair{X}{Y}}{f}$ (rejecting Kuratowski pairs). For each object $A \in \operatorname{Ob}(\mathcal{C})$, the representable presheaf $h_A(X) = \operatorname{Hom}_{\mathcal{C}}(X, A)$ acts contravariantly by precomposition. Given any set-valued presheaf $F\colon \mathcal{C}^{\mathrm{op}} \to \mathbf{Set}$ encoded at the conglomerate level, the Yoneda evaluation map
\[
\Phi_{A,F}\colon \operatorname{Nat}(h_A, F) \longrightarrow F(A), \qquad \Phi_{A,F}(\eta) = \eta_A(1_A),
\]
is a bijection with inverse $\Psi_{A,F}(u)_X(f) = F(f)(u)$ for all $f \in \operatorname{Hom}_{\mathcal{C}}(X, A)$. Specializing to $F = h_B$ shows that the Yoneda embedding $Y(A) = h_A$ is full and faithful: $\operatorname{Nat}(h_A, h_B) \cong \operatorname{Hom}_{\mathcal{C}}(A, B)$. Furthermore, along a functor $K\colon \mathcal{C} \to \mathcal{D}$ and diagram $F\colon \mathcal{C} \to \mathcal{E}$ with Dunaev-tagged graphs $\pair{c}{K(c)}$ and $\pair{c}{F(c)}$, the pointwise left and right Kan extensions are given by the comma-category colimit and limit formulas
\[
(\operatorname{Lan}_K F)(d) \cong \varinjlim_{(K \downarrow d)} F \circ \pi, \qquad
(\operatorname{Ran}_K F)(d) \cong \varprojlim_{(d \downarrow K)} F \circ \pi,
\]
satisfying the universal adjunction identities $\operatorname{Nat}(\operatorname{Lan}_K F, G) \cong \operatorname{Nat}(F, G \circ K)$ and $\operatorname{Nat}(G \circ K, F) \cong \operatorname{Nat}(G, \operatorname{Ran}_K F)$.

\subsection{Elementary Topoi, Subobject Classifiers, and Lawvere-Tierney Topologies}
\label{subsec:topos_subobject_classifier_lawvere_tierney}

An elementary topos $\mathcal{E}$ in \texttt{MKTD1-theory} and \texttt{MKTD2-theory} is a finitely complete, cartesian closed category equipped with a subobject classifier $\top\colon 1 \to \Omega$ such that for every monomorphism $m\colon S \rightarrowtail X$ there exists a unique characteristic morphism $\chi_m\colon X \to \Omega$ making the diagram
\[
\begin{aligned}
S &\longrightarrow 1, \qquad m\colon S \rightarrowtail X, \qquad \chi_m \circ m = \top \circ {!_S}, \\
S &\cong \cc{\pair{x}{u} \in X \times 1}{\chi_m(x) = \top(u)}
\end{aligned}
\]
a pullback square over Dunaev Tagged Ordered Pairs $\pair{x}{u}$. A \emph{Lawvere-Tierney modal topology} on $\mathcal{E}$ is an endomorphism $j\colon \Omega \to \Omega$ satisfying three equational axioms:
\begin{enumerate}
  \item \textbf{Truth Preservation}: $j \circ \top = \top$.
  \item \textbf{Idempotence}: $j \circ j = j$.
  \item \textbf{Meet Preservation}: $j \circ \wedge = \wedge \circ (j \times j)$.
\end{enumerate}
Each Lawvere-Tierney topology $j$ induces a closure operator $S \mapsto \overline{S}^j$ on $\operatorname{Sub}_{\mathcal{E}}(X)$ with characteristic morphism $j \circ \chi_m$, and a subobject $m\colon S \rightarrowtail X$ is $j$-closed precisely when $j \circ \chi_m = \chi_m$.

\subsection{Monads, Eilenberg-Moore and Kleisli Categories, and Beck Monadicity}
\label{subsec:monads_eilenberg_moore_kleisli_beck}

In \texttt{MKTD1-theory} and \texttt{MKTD2-theory}, a monad on a category $\mathcal{C}$ is a triad $\pair{T}{\pair{\eta}{\mu}}$ encoded via Dunaev Tagged Ordered Pairs (rejecting Kuratowski pairs), where $T\colon \mathcal{C} \to \mathcal{C}$ is an endofunctor and $\eta\colon 1_{\mathcal{C}} \Rightarrow T$ (unit) and $\mu\colon T^2 \Rightarrow T$ (multiplication) are natural transformations satisfying the unit and associativity coherence identities:
\[
\mu \circ T\eta = 1_T = \mu \circ \eta T, \qquad \mu \circ T\mu = \mu \circ \mu T.
\]
Each monad $(T, \eta, \mu)$ canonically determines two extremal adjunctions in the category of adjunctions generating $T$:
\begin{enumerate}
  \item \textbf{Eilenberg-Moore Category of $T$-Algebras ($\mathcal{C}^T$, Terminal Resolution)}: Objects are Dunaev-tagged pairs $\pair{A}{h}$ where $h\colon T(A) \to A$ satisfies $h \circ \eta_A = 1_A$ and $h \circ T(h) = h \circ \mu_A$; morphisms $f\colon \pair{A}{h} \to \pair{B}{k}$ satisfy $f \circ h = k \circ T(f)$.
  \item \textbf{Kleisli Category of Free $T$-Algebras ($\mathcal{C}_T$, Initial Resolution)}: Objects are the objects of $\mathcal{C}$, morphisms $f\colon X \rightsquigarrow Y$ in $\mathcal{C}_T$ are morphisms $f\colon X \to T(Y)$ in $\mathcal{C}$ encoded as $\pair{\pair{X}{Y}}{f}$, and Kleisli composition is given by $g \odot f = \mu_Z \circ T(g) \circ f$ with identity $\eta_X$.
\end{enumerate}
Given any adjunction $F \dashv U\colon \mathcal{D} \to \mathcal{C}$ inducing $T = UF$, the canonical comparison functor $K\colon \mathcal{D} \to \mathcal{C}^T$ sends $D \mapsto \pair{U(D)}{U(\varepsilon_D)}$. By Beck's Monadicity Theorem, $K$ is an equivalence of categories if and only if $U$ is conservative (reflects isomorphisms) and $\mathcal{D}$ has and $U$ preserves coequalizers of all $U$-split parallel pairs.

\subsection{Cartesian Closed Categories, Exponential Objects, and Curry-Howard-Lambek Equivalence}
\label{subsec:ccc_exponentials_curry_howard_lambek}

In \texttt{MKTD1-theory} and \texttt{MKTD2-theory}, a category $\mathcal{C}$ with finite products (terminal object $1$ and binary product $A \times B$ encoded via Dunaev Tagged Ordered Pairs $\pair{a}{b}$, rejecting Kuratowski pairs) is \emph{Cartesian closed} if for every pair of objects $A, B \in \operatorname{Ob}(\mathcal{C})$ there exists an \emph{exponential object} $B^A$ together with a canonical \emph{evaluation morphism} $\operatorname{ev}_{A,B}\colon B^A \times A \to B$ satisfying the universal property of exponentiation: for every morphism $g\colon C \times A \to B$, there exists a unique \emph{exponential transpose} (curried morphism) $\lambda(g)\colon C \to B^A$ such that
\[
\operatorname{ev}_{A,B} \circ (\lambda(g) \times 1_A) = g.
\]
Conversely, given $h\colon C \to B^A$, its \emph{uncurried morphism} is $\operatorname{uncurry}(h) = \operatorname{ev}_{A,B} \circ (h \times 1_A)\colon C \times A \to B$, establishing a natural bijection
\[
\operatorname{Hom}_{\mathcal{C}}(C \times A, B) \cong \operatorname{Hom}_{\mathcal{C}}(C, B^A)
\]
whose two-sided inverse equations mirror the $\beta$- and $\eta$-conversion laws of the simply typed $\lambda$-calculus:
\begin{enumerate}
  \item \textbf{Categorical $\beta$-Reduction Law}: $\operatorname{uncurry}(\lambda(g)) = \operatorname{ev}_{A,B} \circ (\lambda(g) \times 1_A) = g$.
  \item \textbf{Categorical $\eta$-Expansion Law}: $\lambda(\operatorname{uncurry}(h)) = \lambda(\operatorname{ev}_{A,B} \circ (h \times 1_A)) = h$.
\end{enumerate}
Under the Curry-Howard-Lambek correspondence, Cartesian closed categories provide the sound and complete categorical semantics for intuitionistic propositional logic (`MKTD2-framework` formulas-as-types) and the simply typed $\lambda$-calculus, interpreting contexts $\Gamma$ as finite Dunaev-tagged products, implication $A \supset B$ as the exponential object $B^A$, modus ponens as $\operatorname{ev}_{A,B}$, and deduction-theorem abstraction as currying $\lambda(-)$.

\subsection{Grothendieck Fibrations, Cartesian Morphisms, Cleavages, and the Grothendieck Construction}
\label{subsec:grothendieck_fibrations_cleavages_construction}

In \texttt{MKTD1-theory} and \texttt{MKTD2-theory}, let $p\colon \mathcal{E} \to \mathcal{B}$ be a functor between categories whose object and morphism graphs are encoded via Dunaev Tagged Ordered Pairs (rejecting Kuratowski pairs). For each base object $I \in \operatorname{Ob}(\mathcal{B})$, the \emph{fiber category} $\mathcal{E}_I = p^{-1}(I)$ consists of objects $X \in \operatorname{Ob}(\mathcal{E})$ with $p(X) = I$ and vertical morphisms $v\colon X \to X'$ satisfying $p(v) = 1_I$. A morphism $\varphi\colon X \to Y$ in $\mathcal{E}$ lying over $u = p(\varphi)\colon I \to J$ in $\mathcal{B}$ is \emph{$p$-Cartesian} if for every object $Z \in \operatorname{Ob}(\mathcal{E})$, every morphism $\psi\colon Z \to Y$ in $\mathcal{E}$, and every base factorization $p(\psi) = u \circ w$ with $w\colon p(Z) \to I$, there exists a unique mediating morphism $\theta\colon Z \to X$ over $w$ ($p(\theta) = w$) such that $\varphi \circ \theta = \psi$:
\[
\forall Z\,\forall \psi\colon Z \to Y\,\forall w\colon p(Z) \to I\;\bigl(p(\psi) = u \circ w \implies \exists!\theta\colon Z \to X\;(p(\theta) = w \wedge \varphi \circ \theta = \psi)\bigr).
\]
Any two $p$-Cartesian lifts of $u\colon I \to J$ with codomain $Y \in \mathcal{E}_J$ are canonically isomorphic via a unique vertical isomorphism in $\mathcal{E}_I$. The functor $p\colon \mathcal{E} \to \mathcal{B}$ is a \emph{Grothendieck fibration} (or \emph{fibered category}) if every base morphism $u\colon I \to J$ and codomain object $Y \in \mathcal{E}_J$ admit a $p$-Cartesian lift $\operatorname{Cart}(u, Y)\colon u^*Y \to Y$; a chosen class-function assignment $(u, Y) \mapsto \operatorname{Cart}(u, Y)$ of such lifts is a \emph{cleavage}, which is \emph{split} when $1_J^* Y = Y$ and $(u \circ v)^* Y = v^*(u^* Y)$ hold on the nose. Given an indexed category (contravariant pseudofunctor) $\Phi\colon \mathcal{B}^{\mathrm{op}} \to \mathbf{Cat}$, the \emph{Grothendieck construction} builds the total category $\int_{\mathcal{B}} \Phi$ whose objects are Dunaev-tagged pairs $\pair{I}{X}$ with $I \in \operatorname{Ob}(\mathcal{B})$ and $X \in \operatorname{Ob}(\Phi(I))$, and whose morphisms $\pair{u}{f}\colon \pair{I}{X} \to \pair{J}{Y}$ consist of $u\colon I \to J$ in $\mathcal{B}$ together with $f\colon X \to \Phi(u)(Y)$ in $\Phi(I)$, establishing the fundamental 2-categorical equivalence between cloven Grothendieck fibrations over $\mathcal{B}$ and indexed categories on $\mathcal{B}^{\mathrm{op}}$.

\subsection{Abelian Categories, Exact Sequences, and the Snake Lemma Connecting Homomorphism}
\label{subsec:abelian_categories_exact_sequences_snake_lemma}

In \texttt{MKTD1-theory} and \texttt{MKTD2-theory}, an \emph{Abelian category} $\mathcal{A}$ (whose objects and morphism graphs are encoded via Dunaev Tagged Ordered Pairs, rejecting Kuratowski pairs) is an additive category equipped with a zero object $0$, binary biproducts $A \oplus B$, kernels $\ker(f) = \operatorname{eq}(f, 0)$, and cokernels $\operatorname{coker}(f) = \operatorname{coeq}(f, 0)$ for all morphisms $f\colon A \to B$, such that every monomorphism is normal (the kernel of its cokernel), every epimorphism is conormal (the cokernel of its kernel), and for every morphism $f\colon A \to B$ the canonical comparison morphism from the coimage to the image is an isomorphism:
\[
\operatorname{coim}(f) = \operatorname{coker}(\ker(f)) \xrightarrow{\;\sim\;} \operatorname{im}(f) = \ker(\operatorname{coker}(f)).
\]
A composable pair $A \xrightarrow{f} B \xrightarrow{g} C$ is \emph{exact at $B$} when $\operatorname{im}(f) = \ker(g)$, and a sequence $0 \to A \xrightarrow{f} B \xrightarrow{g} C \to 0$ is a \emph{short exact sequence} when $f$ is a monomorphism ($\ker(f) = 0$), $g$ is an epimorphism ($\operatorname{coker}(g) = 0$), and $\operatorname{im}(f) = \ker(g)$. Given any commutative diagram of two short exact rows with vertical morphisms $(\alpha, \beta, \gamma)$, the \emph{Snake Lemma} constructs the natural connecting homomorphism $\delta\colon \ker(\gamma) \to \operatorname{coker}(\alpha)$ via the canonical zig-zag diagram chase $c \mapsto b \mapsto \beta(b) \mapsto a' \mapsto [a'] \in \operatorname{coker}(\alpha)$, yielding a six-term exact sequence
\[
\ker(\alpha) \longrightarrow \ker(\beta) \longrightarrow \ker(\gamma) \xrightarrow{\;\delta\;} \operatorname{coker}(\alpha) \longrightarrow \operatorname{coker}(\beta) \longrightarrow \operatorname{coker}(\gamma).
\]

\subsection{Chain Complexes, Projective and Injective Resolutions, and Derived Functors $\mathrm{Ext}^n$ and $\mathrm{Tor}_n$}
\label{subsec:chain_complexes_resolutions_derived_functors_ext_tor}

In \texttt{MKTD1-theory} and \texttt{MKTD2-theory}, let $\mathcal{A}$ be an Abelian category whose graded objects and differential graphs are represented via Dunaev Tagged Ordered Pairs (rejecting Kuratowski pairs). A \emph{chain complex} $(C_\bullet, d_\bullet)$ in $\mathcal{A}$ consists of a $\mathbb{Z}$-indexed family of objects $C_n \in \operatorname{Ob}(\mathcal{A})$ and boundary differentials $d_n\colon C_n \to C_{n-1}$ satisfying the fundamental nilpotency identity $d_{n-1} \circ d_n = 0$ (equivalently $\operatorname{im}(d_{n+1}) \subseteq \ker(d_n)$), with $n$-th homology object $H_n(C_\bullet) = \ker(d_n)/\operatorname{im}(d_{n+1})$. Two chain maps $f_\bullet, g_\bullet\colon C_\bullet \to D_\bullet$ are \emph{chain homotopic} ($f_\bullet \simeq g_\bullet$) via a degree-$+1$ family $s_n\colon C_n \to D_{n+1}$ when $f_n - g_n = d^D_{n+1} \circ s_n + s_{n-1} \circ d^C_n$, inducing identical morphisms $H_n(f_\bullet) = H_n(g_\bullet)$ on homology. When $\mathcal{A}$ has enough projectives (respectively enough injectives), every object $A \in \operatorname{Ob}(\mathcal{A})$ admits a \emph{projective resolution} $P_\bullet \xrightarrow{\varepsilon} A \to 0$ (respectively an \emph{injective resolution} $0 \to A \xrightarrow{\eta} I^\bullet$) unique up to chain homotopy equivalence by the Comparison Theorem; applying the homological and cohomological functors yields the left derived torsion functors and right derived extension functors
\[
\operatorname{Tor}_n^{\mathcal{A}}(A, B) = H_n(P_\bullet \otimes B), \qquad
\operatorname{Ext}^n_{\mathcal{A}}(A, B) = H^n(\operatorname{Hom}_{\mathcal{A}}(P_\bullet, B)) \cong H^n(\operatorname{Hom}_{\mathcal{A}}(A, I^\bullet)),
\]
and every short exact sequence $0 \to A \to B \to C \to 0$ of chain complexes induces a natural \emph{long exact homology/cohomology sequence} via the Snake Lemma connecting homomorphisms $\partial_n\colon H_n(C_\bullet) \to H_{n-1}(A_\bullet)$.

\subsection{Triangulated Categories, Distinguished Triangles, the Octahedral Axiom, and Derived Categories}
\label{subsec:triangulated_categories_octahedral_derived_categories}

In \texttt{MKTD1-theory} and \texttt{MKTD2-theory}, a \emph{triangulated category} $(\mathcal{D}, [1], \Delta)$ is an additive category $\mathcal{D}$ (whose objects, shift functors, and roof fractions are encoded via Dunaev Tagged Ordered Pairs, rejecting Kuratowski pairs) equipped with an additive automorphism $[1]\colon \mathcal{D} \xrightarrow{\;\sim\;} \mathcal{D}$ (the suspension or translation functor $X \mapsto X[1]$) and a class $\Delta$ of \emph{distinguished triangles} $X \xrightarrow{u} Y \xrightarrow{v} Z \xrightarrow{w} X[1]$ satisfying the four Verdier axioms \textup{(TR1)}--\textup{(TR4)}:
\begin{enumerate}
\item \textup{(TR1)} $\Delta$ is closed under isomorphism, contains $X \xrightarrow{1_X} X \to 0 \to X[1]$ for every $X$, and embeds every morphism $u\colon X \to Y$ into a distinguished triangle via the mapping cone $\operatorname{Cone}(u)$;
\item \textup{(TR2)} \emph{Rotation}: $X \xrightarrow{u} Y \xrightarrow{v} Z \xrightarrow{w} X[1]$ is distinguished if and only if $Y \xrightarrow{v} Z \xrightarrow{w} X[1] \xrightarrow{-u[1]} Y[1]$ is distinguished;
\item \textup{(TR3)} \emph{Morphism extension}: any commutative square between the first two arrows of two distinguished triangles extends to a morphism of triangles;
\item \textup{(TR4)} \emph{Octahedral Axiom}: given composable morphisms $X \xrightarrow{u} Y \xrightarrow{v} Z$ and distinguished triangles on $u$, $v$, and $v \circ u$ with third objects $Z'$, $X'$, and $Y'$, there exists a distinguished triangle $Z' \to Y' \to X' \to Z'[1]$ making the octahedral diagram commute.
\end{enumerate}
Given an Abelian category $\mathcal{A}$ and its homotopy category of chain complexes $\mathbf{K}(\mathcal{A})$, the class $S$ of \emph{quasi-isomorphisms} (chain maps inducing isomorphisms on all homology objects $H_n$) forms a localizing multiplicative system in $\mathbf{K}(\mathcal{A})$ admitting a Gabriel--Zisman calculus of left and right fractions (where morphisms in the localization are represented by Dunaev-tagged roofs $X_\bullet \xleftarrow{s} W_\bullet \xrightarrow{f} Y_\bullet$ with $s \in S$), yielding the \emph{derived category} $\mathbf{D}(\mathcal{A}) = \mathbf{K}(\mathcal{A})[S^{-1}]$ together with the universal localization functor $Q\colon \mathbf{K}(\mathcal{A}) \to \mathbf{D}(\mathcal{A})$.

\subsection{Quillen Model Categories, Weak Factorization Systems, and the Homotopy Category}
\label{subsec:quillen_model_categories_lifting_homotopy_category}

In \texttt{MKTD1-theory} and \texttt{MKTD2-theory}, a \emph{Quillen model category} $(\mathcal{M}, \mathcal{W}, \mathcal{C}, \mathcal{F})$ is a complete and cocomplete category $\mathcal{M}$ (whose morphism classes, commutative lifting squares, and cylinder/path factorizations are encoded via Dunaev Tagged Ordered Pairs, rejecting Kuratowski pairs) equipped with three distinguished classes of morphisms---weak equivalences $\mathcal{W}$, cofibrations $\mathcal{C}$, and fibrations $\mathcal{F}$ (with acyclic cofibrations $\mathcal{C} \cap \mathcal{W}$ and acyclic fibrations $\mathcal{F} \cap \mathcal{W}$)---satisfying the five Quillen axioms \textup{(MC1)}--\textup{(MC5)}:
\begin{enumerate}
\item \textup{(MC1)} \emph{Limits and colimits}: $\mathcal{M}$ has all finite limits and colimits (including initial object $\emptyset$ and terminal object $1$);
\item \textup{(MC2)} \emph{Two-out-of-three property}: for composable morphisms $f$ and $g$, if any two of $f$, $g$, and $g \circ f$ belong to $\mathcal{W}$, then so does the third;
\item \textup{(MC3)} \emph{Retracts}: each of $\mathcal{W}$, $\mathcal{C}$, and $\mathcal{F}$ is closed under retracts in the arrow category;
\item \textup{(MC4)} \emph{Lifting ($i \boxslash p$)}: for every commutative square with left vertical arrow $i\colon A \to B$ in $\mathcal{C}$ and right vertical arrow $p\colon X \to Y$ in $\mathcal{F}$, a diagonal lift $h\colon B \to X$ satisfying $h \circ i = f$ and $p \circ h = g$ exists whenever either $i \in \mathcal{C} \cap \mathcal{W}$ or $p \in \mathcal{F} \cap \mathcal{W}$;
\item \textup{(MC5)} \emph{Weak factorization systems}: every morphism $f\colon X \to Y$ factors both as $f = p \circ i$ with $i \in \mathcal{C}$ and $p \in \mathcal{F} \cap \mathcal{W}$, and as $f = q \circ j$ with $j \in \mathcal{C} \cap \mathcal{W}$ and $q \in \mathcal{F}$.
\end{enumerate}
An object $X \in \operatorname{Ob}(\mathcal{M})$ is \emph{cofibrant} if $\emptyset \to X$ is in $\mathcal{C}$, \emph{fibrant} if $X \to 1$ is in $\mathcal{F}$, and \emph{bifibrant} if both hold. By Quillen's Whitehead Theorem, the localization $\operatorname{Ho}(\mathcal{M}) = \mathcal{M}[\mathcal{W}^{-1}]$ is canonically equivalent to the quotient of the full subcategory $\mathcal{M}_{cf}$ of bifibrant objects modulo cylinder/path homotopy $\sim$, and every \emph{Quillen adjunction} $L \dashv R\colon \mathcal{M} \rightleftarrows \mathcal{N}$ (where $L$ preserves cofibrations and acyclic cofibrations, equivalently $R$ preserves fibrations and acyclic fibrations) induces a total derived adjunction $\mathbf{L}L \dashv \mathbf{R}R\colon \operatorname{Ho}(\mathcal{M}) \rightleftarrows \operatorname{Ho}(\mathcal{N})$ which is an equivalence of homotopy categories when $(L, R)$ is a Quillen equivalence.

\subsection{Simplicial Sets, Kan Complexes, Quasi-Categories, and the Homotopy Coherent Nerve}
\label{subsec:simplicial_sets_kan_complexes_quasi_categories}

In \texttt{MKTD1-theory} and \texttt{MKTD2-theory}, let $\boldsymbol{\Delta}$ denote the simplex category of non-empty finite linearly ordered ordinals $[n] = \{0 < 1 < \dots < n\}$ ($n \ge 0$) with order-preserving maps generated by coface maps $\delta^i\colon [n-1] \to [n]$ and codegeneracy maps $\sigma^j\colon [n+1] \to [n]$. A \emph{simplicial set} $X_\bullet\colon \boldsymbol{\Delta}^{\mathrm{op}} \to \mathbf{Set}$ (whose graded $n$-simplices $X_n$ and face/degeneracy operator graphs $d_i\colon X_n \to X_{n-1}$, $s_j\colon X_n \to X_{n+1}$ are encoded via Dunaev Tagged Ordered Pairs, rejecting Kuratowski pairs) satisfies the standard \emph{simplicial identities}:
\[
d_i \circ d_j = d_{j-1} \circ d_i \quad (i < j), \qquad
s_i \circ s_j = s_{j+1} \circ s_i \quad (i \le j), \qquad
d_i \circ s_j = \begin{cases}
s_{j-1} \circ d_i & \text{if } i < j, \\
\mathrm{id} & \text{if } i = j \text{ or } i = j+1, \\
s_j \circ d_{i-1} & \text{if } i > j+1.
\end{cases}
\]
For each $n \ge 1$ and $0 \le k \le n$, the \emph{$k$-th horn} $\Lambda^n_k \subseteq \Delta^n$ is the sub-simplicial set of the standard $n$-simplex $\Delta^n = \boldsymbol{\Delta}(-, [n])$ obtained by removing the interior of $\Delta^n$ and the $k$-th face $d_k(\Delta^n)$; the horn $\Lambda^n_k$ is an \emph{inner horn} when $0 < k < n$ and an \emph{outer horn} when $k = 0$ or $k = n$. A simplicial set $X_\bullet$ is:
\begin{enumerate}
\item a \emph{quasi-category} (or \emph{$\infty$-category} in the Joyal--Boardman--Vogt model) if every inner horn $\Lambda^n_k \to X_\bullet$ ($0 < k < n$) admits a filler $\Delta^n \to X_\bullet$;
\item a \emph{Kan complex} (or \emph{$\infty$-groupoid}, the fibrant objects of the Quillen--Kan model structure) if every horn $\Lambda^n_k \to X_\bullet$ ($0 \le k \le n$, including outer horns) admits a filler $\Delta^n \to X_\bullet$, so every Kan complex is a fortiori a quasi-category;
\item the \emph{categorical nerve} $\operatorname{N}(\mathcal{C})_\bullet$ of a small 1-category $\mathcal{C}$ (whose $n$-simplices are composable $n$-tuples of morphisms $C_0 \to C_1 \to \dots \to C_n$) if and only if every inner horn $\Lambda^n_k \to \operatorname{N}(\mathcal{C})_\bullet$ ($0 < k < n$) admits a \emph{unique} filler, while Cordier's \emph{homotopy coherent nerve} $\operatorname{N}_{\mathrm{hc}}(\mathcal{C}_\Delta)$ extends this construction to simplicially enriched categories $\mathcal{C}_\Delta$ as the right Quillen adjoint to rigidification $\mathfrak{C} \dashv \operatorname{N}_{\mathrm{hc}}$.
\end{enumerate}

\subsection{Grothendieck Sites, Covering Sieves, Sheaf Gluing, and Plus-Construction Sheafification}
\label{subsec:grothendieck_sites_sieves_sheaf_gluing_sheafification}

In \texttt{MKTD1-theory} and \texttt{MKTD2-theory}, let $\mathcal{C}$ be a category whose objects, morphism graphs, and sieve subfunctors are encoded via Dunaev Tagged Ordered Pairs (rejecting Kuratowski pairs). For an object $U \in \operatorname{Ob}(\mathcal{C})$, a \emph{sieve on $U$} is a right ideal of morphisms with codomain $U$ (equivalently a subfunctor $S \hookrightarrow h_U = \operatorname{Hom}_{\mathcal{C}}(-, U)$ closed under precomposition: $f \in S(V) \implies f \circ g \in S(W)$ for all $g\colon W \to V$). Given a morphism $h\colon V \to U$, the \emph{pullback sieve} $h^*(S)$ on $V$ is defined by $h^*(S)(W) = \{g\colon W \to V \mid h \circ g \in S(W)\}$. A \emph{Grothendieck topology} $J$ on $\mathcal{C}$ assigns to each object $U$ a collection $J(U)$ of \emph{covering sieves} on $U$ satisfying:
\begin{enumerate}
\item \emph{Maximal sieve axiom}: the maximal sieve $t_U = \operatorname{Hom}_{\mathcal{C}}(-, U)$ belongs to $J(U)$;
\item \emph{Pullback stability}: if $S \in J(U)$ and $h\colon V \to U$, then $h^*(S) \in J(V)$;
\item \emph{Local character (transitivity)}: if $S \in J(U)$ and $R$ is any sieve on $U$ such that $h^*(R) \in J(V)$ for all $h \in S(V)$, then $R \in J(U)$.
\end{enumerate}
The pair $(\mathcal{C}, J)$ is called a \emph{Grothendieck site}. Given a presheaf $P\colon \mathcal{C}^{\mathrm{op}} \to \mathbf{Set}$ and a sieve $S \in J(U)$, a \emph{matching family} for $P$ on $S$ is a natural transformation $x\colon S \Rightarrow P$ (a family $\{x_f \in P(\operatorname{dom}(f))\}_{f \in S}$ satisfying $P(g)(x_f) = x_{f \circ g}$), and an \emph{amalgamation} is an element $s \in P(U)$ such that $P(f)(s) = x_f$ for all $f \in S$. The presheaf $P$ is a \emph{$J$-separated presheaf} if every matching family has at most one amalgamation, and a \emph{$J$-sheaf} if every matching family on every covering sieve $S \in J(U)$ has a \emph{unique} amalgamation ($\operatorname{Hom}(h_U, P) \xrightarrow{\sim} \operatorname{Hom}(S, P)$). Grothendieck's \emph{plus-construction} $P^+(U) = \varinjlim_{S \in J(U)^{\mathrm{op}}} \operatorname{Match}(S, P)$ maps any presheaf $P$ to a separated presheaf $P^+$ and maps any separated presheaf to a sheaf, so the iterated functor $\mathbf{a}(P) = P^{++}\colon \mathbf{PSh}(\mathcal{C}) \to \mathbf{Sh}(\mathcal{C}, J)$ is the exact left adjoint (\emph{sheafification}) to the full inclusion $\mathbf{Sh}(\mathcal{C}, J) \hookrightarrow \mathbf{PSh}(\mathcal{C})$, yielding the Grothendieck topos $\mathbf{Sh}(\mathcal{C}, J)$.

\subsection{Spectral Sequences, Filtered Complexes, Exact Couples, and Grothendieck Convergence}
\label{subsec:spectral_sequences_filtered_complexes_grothendieck}

In \texttt{MKTD1-theory} and \texttt{MKTD2-theory}, let $\mathcal{A}$ be an Abelian category whose bigraded objects, differentials, and filtration stages are encoded via Dunaev Tagged Ordered Pairs (rejecting Kuratowski pairs). A \emph{cohomological spectral sequence} in $\mathcal{A}$ (starting at page $r_0 \ge 0$) consists of a family of bigraded objects $\{E_r^{p,q}\}_{p,q \in \mathbb{Z},\, r \ge r_0}$ together with differentials $d_r^{p,q}\colon E_r^{p,q} \to E_r^{p+r, q-r+1}$ of bidegree $(r, 1-r)$ satisfying $d_r \circ d_r = 0$, together with canonical page-transition isomorphisms
\[
E_{r+1}^{p,q} \;\cong\; H^{p,q}(E_r^{\bullet,\bullet}, d_r) \;=\; \frac{\ker\bigl(d_r^{p,q}\colon E_r^{p,q} \to E_r^{p+r, q-r+1}\bigr)}{\operatorname{im}\bigl(d_r^{p-r, q+r-1}\colon E_r^{p-r, q+r-1} \to E_r^{p,q}\bigr)}.
\]
Equivalently, any \emph{exact couple} $(D, E, i, j, k)$ in $\mathcal{A}$---an exact triangle $D \xrightarrow{i} D \xrightarrow{j} E \xrightarrow{k} D$---generates a derived exact couple $(D', E', i', j', k')$ with $D' = \operatorname{im}(i)$, $E' = H(E, j \circ k)$, and differential $d = j \circ k\colon E \to E$ satisfying $d \circ d = j \circ (k \circ j) \circ k = 0$. Given a bounded decreasingly \emph{filtered cochain complex} $\cdots \supseteq F^p C^\bullet \supseteq F^{p+1} C^\bullet \supseteq \cdots$ (exhaustive and Hausdorff), the spectral sequence with $E_0^{p,q} = \operatorname{Gr}^p C^{p+q} = F^p C^{p+q}/F^{p+1} C^{p+q}$ and $E_1^{p,q} = H^{p+q}(\operatorname{Gr}^p C^\bullet)$ converges strongly to the filtered cohomology abutment:
\[
E_1^{p,q} \;\Longrightarrow\; H^{p+q}(C^\bullet), \qquad E_\infty^{p,q} \;\cong\; \operatorname{Gr}^p H^{p+q}(C^\bullet) = \frac{F^p H^{p+q}(C^\bullet)}{F^{p+1} H^{p+q}(C^\bullet)}.
\]
In particular, if $F\colon \mathcal{A} \to \mathcal{B}$ and $G\colon \mathcal{B} \to \mathcal{C}$ are left exact functors between Abelian categories with enough injectives such that $F$ sends injective objects of $\mathcal{A}$ to $G$-acyclic objects of $\mathcal{B}$ ($(\mathbf{R}^q G)(F(I)) = 0$ for all $q > 0$), the \emph{Grothendieck composite-functor spectral sequence} satisfies
\[
E_2^{p,q} = (\mathbf{R}^p G)\bigl((\mathbf{R}^q F)(A)\bigr) \;\Longrightarrow\; \mathbf{R}^{p+q}(G \circ F)(A).
\]

\subsection{Symmetric Monoidal Closed Categories, Enriched Yoneda, and Day Convolution}
\label{subsec:symmetric_monoidal_closed_enriched_yoneda_day_convolution}

In \texttt{MKTD1-theory} and \texttt{MKTD2-theory}, let $\mathcal{V} = (\mathcal{V}_0, \otimes, I, \alpha, \lambda, \rho, \sigma)$ be a \emph{symmetric monoidal category} whose objects, tensor product graphs, coherence natural isomorphisms, and enriched hom-objects are encoded via Dunaev Tagged Ordered Pairs (rejecting Kuratowski pairs). By \emph{Mac Lane's Coherence Theorem}, every formal diagram built from the associator $\alpha_{A,B,C}\colon (A \otimes B) \otimes C \xrightarrow{\sim} A \otimes (B \otimes C)$, left and right unitors $\lambda_A\colon I \otimes A \xrightarrow{\sim} A$ and $\rho_A\colon A \otimes I \xrightarrow{\sim} A$, and involutive braiding $\sigma_{A,B}\colon A \otimes B \xrightarrow{\sim} B \otimes A$ ($\sigma_{B,A} \circ \sigma_{A,B} = \operatorname{id}_{A \otimes B}$) commutes whenever the \emph{pentagon identity}
\[
(\operatorname{id}_A \otimes \alpha_{B,C,D}) \circ \alpha_{A,B \otimes C,D} \circ (\alpha_{A,B,C} \otimes \operatorname{id}_D) \;=\; \alpha_{A,B,C \otimes D} \circ \alpha_{A \otimes B,C,D}
\]
and the \emph{triangle identity} $(\operatorname{id}_A \otimes \lambda_B) \circ \alpha_{A,I,B} = \rho_A \otimes \operatorname{id}_B$ (together with hexagon coherence) hold. Furthermore, $\mathcal{V}$ is \emph{symmetric monoidal closed} if for every object $B \in \operatorname{Ob}(\mathcal{V})$, the endofunctor $(-) \otimes B\colon \mathcal{V} \to \mathcal{V}$ admits a right adjoint $[B, -]\colon \mathcal{V} \to \mathcal{V}$ (\emph{internal hom}), yielding natural tensor-hom adjunction isomorphisms
\[
\operatorname{Hom}_{\mathcal{V}}(A \otimes B, C) \;\cong\; \operatorname{Hom}_{\mathcal{V}}\bigl(A, [B, C]\bigr), \qquad [A \otimes B, C] \;\cong\; \bigl[A, [B, C]\bigr].
\]
When $\mathcal{V}$ is complete and cocomplete symmetric monoidal closed, a \emph{$\mathcal{V}$-enriched category} $\underline{\mathcal{C}}$ consists of a class of objects $\operatorname{Ob}(\underline{\mathcal{C}})$, hom-objects $\underline{\mathcal{C}}(X, Y) \in \operatorname{Ob}(\mathcal{V})$, composition morphisms $M_{X,Y,Z}\colon \underline{\mathcal{C}}(Y, Z) \otimes \underline{\mathcal{C}}(X, Y) \to \underline{\mathcal{C}}(X, Z)$, and unit elements $j_X\colon I \to \underline{\mathcal{C}}(X, X)$ satisfying enriched associativity and unit laws. For any $\mathcal{V}$-functor $F\colon \underline{\mathcal{C}} \to \mathcal{V}$, the \emph{enriched Yoneda Lemma} establishes an isomorphism in $\mathcal{V}$ via the enriched end
\[
F(A) \;\cong\; [\underline{\mathcal{C}}, \mathcal{V}]\bigl(\underline{\mathcal{C}}(A, -), F\bigr) \;=\; \int_{X \in \underline{\mathcal{C}}} \bigl[\underline{\mathcal{C}}(A, X), F(X)\bigr].
\]
Finally, if $\underline{\mathcal{C}}$ is a small symmetric monoidal $\mathcal{V}$-category, the enriched functor category $[\underline{\mathcal{C}}^{\mathrm{op}}, \mathcal{V}]$ inherits a symmetric monoidal closed structure via \emph{Day convolution} coend
\[
(F \star_{\mathrm{Day}} G)(Z) \;=\; \int^{X, Y \in \underline{\mathcal{C}}} \underline{\mathcal{C}}(Z, X \otimes_{\mathcal{C}} Y) \otimes F(X) \otimes G(Y),
\]
such that the enriched Yoneda embedding $\mathbf{y}\colon \underline{\mathcal{C}} \hookrightarrow [\underline{\mathcal{C}}^{\mathrm{op}}, \mathcal{V}]$ is strong symmetric monoidal: $\mathbf{y}(A \otimes_{\mathcal{C}} B) \cong \mathbf{y}(A) \star_{\mathrm{Day}} \mathbf{y}(B)$.

\ifdefined\qwertyflag
\undef\qwertyflag
\end{document}
\fi
